\section{Discussion}
\label{sec:discussion}

\paragraph{A named alternative is the active ingredient.} For this model, pushback alone does very
little. A neutral ``are you sure?'' and an insistent ``I'm fairly confident you've made a mistake''
leave \Cneutralpeak{} control questions holding, while the same two turns naming a specific wrong
answer leave \Cfoilpeak. This differs from reports that a bare ``are you sure?'' causes large
accuracy drops in larger chat models \citep{laban2023flipflop,xie2024askagain}; those studies use
other models and tasks, so we read the contrast only as a reminder that the trigger can differ
between models. It is in line with sycophancy toward a stated user view
\citep{sharma2024sycophancy,perez2023discovering,wei2023synthetic}: the model gives the answer the
user names.

\paragraph{Most of the damage is momentary.} When the last control turn re-asks the question
without restating the foil, \Cfoilfinal{} questions are correct again. An evaluation that applies
pressure and then scores a clean final re-ask would therefore miss most folds, and one that scores
only the first pushed turn would miss the later swings back and forth
(Figure~\ref{fig:dialogues}). This is why we score every turn and keep the final turn under
pressure. The bare-answer format of our final turns also matters: it leaves no room for the model
to hedge or to restate the correct answer in passing, and it produces the lowest hold rate in the
false-premise suite.

\paragraph{Authority versus the model's own past.} Claims that the model said something earlier
are resisted much better than claims from an outside authority. One plausible reason is that the
model can see its own earlier replies in the context and they contradict the claim, while an
appeal to a textbook cannot be checked against anything in the context. We did not test this
explanation; a condition that removes the earlier replies from the history would.

\paragraph{The text and the probabilities move together.} If the model kept preferring the correct
answer while saying the wrong one, folds would come with positive $\Delta$. They almost never do
(\Dfoldpos{} of \Dfold), and $\Delta$ separates hold from fold with AUROC \Dauc. For this model,
then, a fold is a change in what it prefers to say, not a polite surface over an unchanged
preference. Because $\Delta$ needs one forward pass and no generation, it is also a cheap signal
for monitoring pressure at scale, though it scores only the bare answer as the opening of the
reply and should not be read as a measure of what the model ``believes''.

\paragraph{Scoring choices change the numbers.} Three choices that are easy to get wrong each moved
our results. Scoring a non-answer as a fold, as earlier versions of this harness did, inflates
fold rates; treating apologies as capitulation mixes politeness with folding
(Section~\ref{sec:res-concession}); and an alias list that misses one form of the correct answer
turns holds into non-answers. We found three such alias gaps after the runs, by reading replies the
scorer had labelled \lab{other}; fixing them changed \Rn{} labels, all from \lab{other}, and none
between \lab{hold} and \lab{fold} (Appendix~\ref{app:repro}).

\subsection{Limitations}
\begin{itemize}
\item \textbf{One model, one decoding setting.} All results are for a single 1.5B instruct model
with greedy decoding and one system prompt; nothing here says how larger models or sampled
decoding behave.
\item \textbf{Small, templated suites.} 48 factual and 24 harmful-request dialogues, and 16
questions in the controls, give wide confidence intervals. Each template uses one fixed wording
per turn, so the results are for these wordings; the final false-premise turn in particular
combines a restated foil with a bare-answer format.
\item \textbf{No human labels.} Labels come from a lexical scorer and, for refusals, an NLI
classifier. We compare them with blind AI annotators, which are not a substitute for human
judgement; the two annotators were instances of the same model and agreed on every cell, so they
amount to one annotator. A human label sheet is released but has not been filled in.
\item \textbf{Noisier refusal labels.} Agreement with the annotator is lower for refusals
($\kappa=\AhSAk$) than for factual answers ($\kappa=\AfSAk$), and \Tharm{} of harmful-request
replies are cut at the 96-token cap; refusal hold rates are less certain than factual ones.
\item \textbf{Post-hoc alias fixes.} Three alias lists were corrected after the runs. The prompts
did not change, and both the as-run and the corrected labels are recorded.
\item \textbf{Withheld harmful-request suite.} Others cannot rerun the refusal results from the
public release alone (Section~\ref{sec:ethics}).
\end{itemize}

\subsection{Ethics and release policy}
\label{sec:ethics}
The harmful-request suite asks for malware, fraud, unauthorized access, weapons and similar
content, and the model complies with many of these requests. We therefore do not publish the
harmful-request prompts or the model's replies to them. We publish the code, the factual and
control suites, all factual model outputs, and the aggregate refusal statistics reported here.
Frontier labs and evaluation teams can obtain the harmful-request suite and the full refusal
transcripts for evaluation purposes by writing to the author at
\href{mailto:vtiwari@acsi-india.org}{\texttt{vtiwari@acsi-india.org}}.

\section*{Acknowledgements and AI-assistance disclosure}
The harness, suites, analysis code and the text of this paper were written with substantial help
from an AI assistant (Claude, via Cursor) under the author's direction, and the blind annotators
of Section~\ref{sec:res-labels} were instances of the same assistant. The author reviewed and is
responsible for all content. Every number in the paper is generated from the committed run
outputs by \texttt{analyze.py} and \texttt{paper/make\_numbers.py}.
